\documentclass[10pt,a4paper]{amsart}
\usepackage[margin=19mm]{geometry}
\usepackage{amsmath,amssymb,mathtools}
\usepackage[scr=rsfs,cal=euler]{mathalfa}
\usepackage{xfrac}
\usepackage[unicode,hidelinks,bookmarks=false]{hyperref}
\newcommand{\ie}{i.e.\ }
\newcommand{\cf}{cf.\ }
\newcommand{\resp}{respectively\ }
\newcommand{\Subsection}{\subsection}
\providecommand{\nobreakdash}{\nobreak\hbox{-}}
\setlength{\emergencystretch}{2em}
\multlinegap0pt
\usepackage{siunitx}
\usepackage{siunitx}
\usepackage{siunitx}
\usepackage{siunitx}
\usepackage{amssymb}
\usepackage{mathtools}
\usepackage{bm}
\usepackage{dsfont}
\usepackage{cancel}
\usepackage{siunitx}
\newtheorem{theorem}{Theorem}
\title{Information Rate Decomposition for Noisy Nanopore Channels with Geometric Duplication}
\author{Brendon McBain, Emanuele Viterbo}
\date{Source: arXiv:2606.06808v2 (CC BY 4.0)}
\begin{document}
\maketitle

\noindent\emph{Source and licence.} Information Rate Decomposition for Noisy Nanopore Channels with Geometric Duplication. Version arXiv:2606.06808v2 (CC BY 4.0), distributed by arXiv under the Creative Commons Attribution 4.0 International licence, http://creativecommons.org/licenses/by/4.0/, as recorded in the arXiv licence metadata for this exact version and shown on the abstract page https://arxiv.org/abs/2606.06808v2. Full licence notice: \texttt{licenses/arxiv-2606.06808-CC-BY-4.0.txt}.\par\smallskip

\noindent\textit{Editorial scope.} Counted unit: the article's theorem thm:jump\_reliability\_lower\_bound (source0.tex lines 945--977). The article's complete original proof of it is delivered with it (source0.tex lines 979--1044, one \texttt{proof} environment of 66 lines). No mathematics is written, repaired or re-derived: the statement and the proof are the article's own lines, in the article's own notation, and no retained line is substituted or re-flowed.  No further source-local context is needed: every cross-reference of every retained line resolves inside the two delivered spans.  Nothing was omitted from any retained span, so the package carries no omission record.  No retained line contains a citation, so the wrapper carries no bibliography and no bibliography entry.  No retained line contains a figure environment, an image-inclusion command or drawing code, so no image is delivered and the package declares no asset dependency.  Editorial formatting only, in addition: an AMS \texttt{amsart} preamble that restates the article's own declarations and macros used by the retained lines, namely the environments \texttt{theorem} on separate counters, together with the standard packages those lines need.  The retained environments print in this document as Theorem 1 in order of appearance, whereas the article prints the same items with the numbering of its own full document.  The complete native source of the article as distributed in the arXiv e-print for this exact version is bundled unchanged as \texttt{sources/202214/source0.tex}.
\begin{theorem}[Jump-reliability lower bound]
\label{thm:jump_reliability_lower_bound}
Let 
\begin{align}
    \overline{R}_{\rm seg}
    =
    \frac{1}{\mu}
    I(K_1;{T_2},Y^{T_2}\mid S_1,S_2).
\end{align}
Then
\begin{align}\label{eq:LB_eq}
    I_T(S;Y)\ge \mu\bigl(I_{\rm ISI}-\overline{R}_{\rm seg}\bigr).
\end{align}
Moreover,
{\begin{align}
    \overline{R}_{\rm seg}
    &=h_b\left(\frac{1}{\mu}\right)- \frac{1}{\ln 2}
    \biggl( 1-\sum_{s',s\in\Omega} q_0(s')P_{S|S^-}(s\mid s')
    \overline{\psi}(s',s)
    \biggr),
\end{align}}
with
{\begin{align}
    \overline{\psi}(s',s)
    &=
    \frac{1}{2\mu\sigma^2}
    \mathbb{E}\!\left[
        \Psi_{2\sigma^2}\bigl((f(s'),f(s)),Y^{T_2}\bigr)
        \,\middle|\,
        S_1=s',S_2=s
    \right].
\end{align}}
\end{theorem}

\begin{proof}
Starting from the chain rule,
\begin{align*}
    I(K^m;Y^{T_m}\mid S^m)
    &= \sum_{\ell=1}^m
    I(K_\ell;Y^{T_m}\mid S^m,K^{\ell-1}) \\
    &= \sum_{\ell=1}^m
    I(K_\ell;Y_{T_{\ell-1}+1}^{T_m}
    \mid S^m,T_{\ell-1}).
\end{align*}
For each $\ell\leq m-1$, we append the neighbouring jump time
$T_{\ell+1}$ to the observations. Once $T_{\ell+1}$ is revealed, the
outputs following $T_{\ell+1}$ are conditionally independent of
$K_\ell$. Therefore,
\begin{align*}
    I(K^m;Y^{T_m}\mid S^m) &\leq
    \sum_{\ell=1}^{m-1}
    I\!\left(
        K_\ell;
        {T_{\ell+1}},Y_{T_{\ell-1}+1}^{T_{\ell+1}}
        \,\middle|\,
        S^m,T_{\ell-1}
    \right) + {H(K_m\mid S_m)} \\
    &=
    (m-1)
    I(K_1;{T_2},Y^{T_2}\mid S_1,S_2)
    +{H(K_m\mid S_m)} \\
    &=
    \mu(m-1)\overline{R}_{\rm seg}
    +{H(K_m\mid S_m)},
\end{align*}
where the second equality follows from stationarity under shifts in the
segment index and conditional independence. It follows that
\begin{align*}
    I_T(K;Y\mid S)
    &\leq
    \lim_{m\rightarrow\infty}
    \left(
        \mu\frac{m-1}{m}\overline{R}_{\rm seg}
        +{\frac{1}{m}H(K_m\mid S_m)}
    \right) =\mu\overline{R}_{\rm seg}.
\end{align*}
Substituting this upper bound on the segmentation term into the rate
decomposition yields \eqref{eq:LB_eq}.

For the expression of $\overline{R}_{\rm seg}$, we have
\begin{align*}
    \overline{R}_{\rm seg}
    &=
    \frac{1}{\mu}
    I(K_1;T_2,Y^{T_2}\mid S_1,S_2) =
    \frac{1}{\mu}
    \left[
        H(K)
        -H(K_1\mid Y^{T_2},S_1,S_2,T_2)
    \right].
\end{align*}
Conditioning further on $(S_1,S_2)=(s',s)$ gives
\begin{align}
    H(K_1\mid Y^{T_2},S_1=s',S_2=s,T_2)
    =
    \frac{\mu}{\ln 2}
    \bigl(1-\overline{\psi}(s',s)\bigr),\label{eq:two-segment-ent}
\end{align}
and averaging over $(S_1,S_2)$ completes the proof.
\end{proof}

\end{document}
